%% mwe_tabelas2_phase3.tex
%% Substitutos do tabularray sob phase-III:
%%   D: tabular + \rowcolor (colortbl via xcolor[table])
%%   E: tabular com coluna p{} (texto que quebra linha)
%%   F: longtable (tabela que atravessa páginas) + \rowcolor
\DocumentMetadata{
  lang        = pt-BR,
  pdfversion  = 2.0,
  pdfstandard = ua-2,
  testphase   = {phase-III, table, firstaid},
}
\documentclass[11pt]{article}
\usepackage[T1]{fontenc}
\usepackage[a4paper, margin=25mm]{geometry}
\usepackage[table]{xcolor}
\usepackage{longtable}
\usepackage{booktabs}
\definecolor{herhighness}{HTML}{422C73}
\colorlet{alarmefundo}{red!10!white}
\setlength{\parindent}{0pt}

\begin{document}

\section*{D --- tabular + rowcolor}

\begin{tabular}{ll}
  \rowcolor{herhighness}
  \textcolor{white}{\textbf{Etapa}} & \textcolor{white}{\textbf{Responsável}} \\
  Definição de pauta & ComCom \\
  \rowcolor{alarmefundo}
  \textbf{Checagem de fatos} & \textbf{Comunicação} \\
  Validação final & Plenário \\
\end{tabular}

\section*{E --- tabular com p\{\} e booktabs}

\begin{tabular}{p{0.35\linewidth}p{0.5\linewidth}}
  \toprule
  \textbf{Etapa} & \textbf{Responsável na Opção B} \\
  \midrule
  Definição de pauta & ComCom mais Comunicação, conforme já previsto
  no Regimento Interno do Conselho \\
  Checagem de fatos & Comunicação \\
  \bottomrule
\end{tabular}

\section*{F --- longtable}

\begin{longtable}{p{0.2\linewidth}p{0.6\linewidth}}
  \rowcolor{herhighness}
  \textcolor{white}{\textbf{Ano}} & \textcolor{white}{\textbf{Observação}} \\
  \endfirsthead
  \rowcolor{herhighness}
  \textcolor{white}{\textbf{Ano}} & \textcolor{white}{\textbf{Observação}} \\
  \endhead
  1981 & Primeira edição da série, com equipe nominal reduzida \\
  1999 & Reestruturação editorial \\
  2020 & Internalização da produção \\
  2026 & Menor equipe nominal da série \\
\end{longtable}

\end{document}
